% 第3章 金利の期間構造(資料3)
\chapter{金利の期間構造}

利回りは投資期間(残存期間)によって異なる値をとる.この「期間と利回りの関係」を
\textbf{金利の期間構造} (term structure of interest rates) という.本章では,
将来のキャッシュフローが確定している場合を考え,安全確実な債券($\approx$ 国債)を
対象として,最終利回り,イールドカーブ,ゼロレート,フォワード・イールドという
一連の概念を学ぶ.

\section*{導入:国債利回りのデータ}

講義では,日本相互証券が公表する10年国債利回りデータ(BB国債価格(引値))を
もとに,各年度末時点の利回りカーブが示された.残存期間(年)を横軸,複利利回り
(\%)を縦軸にとると,例えば次のような特徴が観察される.

\begin{itemize}
\item 2026年3月末:残存期間とともに利回りが上昇する右上がりのカーブで,
10年で2\%を超える水準.
\item 2025年3月末:同じく右上がりで,10年で1.5\%前後.
\item 2024年3月末:右上がりで,10年で0.7\%程度.
\item 2023年2月末:短期はマイナス圏($-0.1\%$前後)から始まり,10年で0.5\%程度.
\item 2022年2月末:短期が$-0.05\%$程度,10年で0.2\%弱.
\item 2021年3月末:短中期が広くマイナス圏で,10年でようやく0.1\%程度.
\item 2020年3月末:7年超までほぼフラット(マイナス圏)という特異な形状.
\end{itemize}

このように,利回りは残存期間ごとに異なり,カーブ全体の水準も形状も時期によって
大きく変化する.また,2000年代(2000〜2007年の各年度末)の国債の金利期間構造
(ゼロレートカーブ)を重ねて描くと,どの年も右上がり(順イールド)であることが
確認できる.本章はこの「期間構造」を記述する道具立てを与える.

\section{最終利回りと債券価格}

\subsection{証券の現在価値とディスカウントファクター}

複数回の確定的な将来キャッシュフローの場合,前章で学んだとおり,
\[
\text{証券の現在価値} = \text{(将来キャッシュフローの現在価値の総和)}
= \sum_i \bigl\{(\text{時刻 } t_i \text{ の将来CF}) \times
(\text{時刻 } t_i \text{ における1円の現在価値})\bigr\}
\]
である.「ある証券を持つこと」と「その証券から生じる個々の将来キャッシュフローを
個々の契約とみなし,すべての契約をまとめたものを持つこと」の経済価値は同じであり,
後者で考える方がわかりやすい.

時刻 $t$ のキャッシュフローを $CF(t)$,時刻 $t$ における1円の現在価値を
$DF(t)$ と書く.キャッシュフローのある時点を $t_i$, $i=1,2,\dots,N$ とすると,
現在価値 $PV$ は
\begin{equation}
PV = \sum_{i=1}^{N} CF(t_i)\, DF(t_i).
\label{eq:pv-df}
\end{equation}

\begin{definition}[ディスカウントファクター]
将来時点 $t$ における1円の現在価値 $DF(t)$ を,\textbf{ディスカウント
ファクター}(ディスカウントファンクション,割引関数)と呼ぶ.
\end{definition}

\subsection{最終利回りの定義}

額面 $F$ 円,満期2年,クーポンレート $C$,年1回払の満期一括償還の利付債を
考える.価格を $P$ として,
\begin{equation}
P = \frac{FC}{1+Y} + \frac{FC}{(1+Y)^2} + \frac{F}{(1+Y)^2}
\end{equation}
という式を満たす $Y$ を\textbf{最終利回り} (Yield To Maturity, YTM) と呼ぶ.

同じ債券でも年2回払であれば,複利の設定(転化回数 $m$)に応じて式が変わり,
\begin{equation}
P = \frac{FC/2}{1+Y/2} + \frac{FC/2}{(1+Y/2)^2} + \frac{FC/2}{(1+Y/2)^3}
+ \frac{FC/2}{(1+Y/2)^4} + \frac{F}{(1+Y/2)^4}
\end{equation}
となる.

\begin{definition}[最終利回り(一般形)]
額面 $F$ 円,満期 $N$ 年,クーポンレート $C$,年 $m$ 回払の満期一括償還の
利付債の価格を $P$ とするとき,
\begin{equation}
P = \frac{F}{(1+Y/m)^{mN}} + \sum_{j=1}^{mN}\frac{FC/m}{(1+Y/m)^{j}}
\label{eq:ytm-def}
\end{equation}
を満たす $Y$ を\textbf{最終利回り} (Yield To Maturity) と呼ぶ.
\end{definition}

前章までの考え方によると,式\eqref{eq:ytm-def}は「満期まで一定の利回り $Y$ を
用いて将来キャッシュフローを割り引いた現価(現在価値)が $P$」と読める.

\begin{keypoint}[債券価格と最終利回りの関係]
債券価格 $P$ と最終利回り $Y$ は,\textbf{逆方向に動く}.一方が上がれば他方は
下がり,一方が下がれば他方は上がる.すなわち関数 $P(Y)$ は単調減少である.
最終利回り $Y$ は,価格 $P$ を利回りで言い換えたものである.
\end{keypoint}

講義資料では次の2つの例が図示された.
\begin{itemize}
\item 5年満期割引国債(額面100円):最終利回りが0\%のとき価格100円,利回りが
上がるにつれて価格は単調に下落し,5\%で約78円となる.
\item 10年満期,クーポンレート3\%,額面100円,満期一括償還の利付債:利回り0\%の
とき価格130円,\textbf{最終利回り $=$ クーポンレートのとき価格はちょうど100円}
(パー)となり,8\%では約66円まで下がる.
\end{itemize}
最終利回りとクーポンレートが等しいとき価格が額面に一致することは,
式\eqref{eq:ytm-def}に $Y=C$ を代入すると等比数列の和が畳み込めて $P=F$ となる
ことから確認できる(まず1期間のケースで確かめるとよい).

\begin{example}[例題3-1:債券の現在価値]
額面100円,満期2年,クーポンレート6\%,半年払い($m=2$)の債券の現在価値を
求めなさい.
\begin{enumerate}
\item[(1)] 複利利回りを5\%(半年複利)とする場合.
\item[(2)] 連続複利利回りを5\%とする場合.
\end{enumerate}

\medskip
\noindent\textbf{解答.}
キャッシュフローは,0.5年,1年,1.5年後に各3円($=6\%\times 100\times 0.5$),
2年後に103円である.

(1) 半年複利のディスカウントファクターは $DF(t) = (1+0.05/2)^{-2t}$ であるから,
\[
PV = \sum_{i=1}^{4} CF(0.5i)\,DF(0.5i)
= \frac{3}{1+\frac{0.05}{2}} + \frac{3}{(1+\frac{0.05}{2})^2}
+ \frac{3}{(1+\frac{0.05}{2})^3} + \frac{103}{(1+\frac{0.05}{2})^4}
\approx 101.88.
\]

(2) 連続複利のディスカウントファクターは $DF(t) = e^{-0.05t}$ であるから,
\[
PV = 3e^{-0.05\times 0.5} + 3e^{-0.05\times 1} + 3e^{-0.05\times 1.5}
+ 103e^{-0.05\times 2} \approx 101.76.
\]

計算の内訳は次表のとおりである.
\begin{center}
\small
\begin{tabular}{cccc@{\qquad}cccc}
\toprule
\multicolumn{4}{c}{半年複利(利回り5.0\%)} & \multicolumn{4}{c}{連続複利(利回り5.0\%)}\\
\cmidrule(r){1-4}\cmidrule(l){5-8}
年 & CF & DF & 価値 & 年 & CF & DF & 価値\\
\midrule
0.5 & 3 & 0.97561 & 2.92683 & 0.5 & 3 & 0.97531 & 2.92593\\
1.0 & 3 & 0.95181 & 2.85544 & 1.0 & 3 & 0.95123 & 2.85369\\
1.5 & 3 & 0.92860 & 2.78580 & 1.5 & 3 & 0.92774 & 2.78323\\
2.0 & 103 & 0.90595 & 93.31292 & 2.0 & 103 & 0.90484 & 93.19825\\
\midrule
\multicolumn{3}{r}{現在価値} & 101.88099 & \multicolumn{3}{r}{現在価値} & 101.7611\\
\bottomrule
\end{tabular}
\end{center}
\end{example}

\subsection{価格から最終利回りを求める}

定義式\eqref{eq:ytm-def}に沿って,価格 $P$ から最終利回り $Y$ を算出することを
考える.一般に,これは $mN$ 次方程式を解くことになるため,$mN \geq 2$ では
$Y$ を求めるには数値計算が必要である.

注意点として,(高次方程式なので)不適切な解も存在しえるため,結果をチェックする
必要がある.とはいえ,$Y$ を求める関数はエクセルにも実装されていて,一般的に
使われている.

\section{イールドカーブ}

\begin{definition}[イールドカーブ(利回り曲線)]
個々の債券について,残存期間を $x$ 座標,利回りを $y$ 座標として平面上に
描いた線(カーブ)のこと.利回りの期間構造(利回り曲線)ともいう.
\end{definition}

利回りにはいろいろな種類のものがあるが,債券投資の実務では「最終利回り」が
よく使われる.

カーブの概形については次の用語がある.
\begin{itemize}
\item \textbf{順イールド} (normal yield):右上がりのカーブ.
\item \textbf{逆イールド} (inverted yield):右下がりのカーブ.
\end{itemize}
実際には順イールドになることが多い.このほか,中間の期間が盛り上がった
humped shape と呼ばれる形状もある.

先に見た2026年3月末の10年国債利回りカーブは典型的な順イールドである.一方,
2020年3月末のカーブは,7年超までほぼフラットでマイナス圏という形状であり,
「順イールド?」と問いたくなるような特異な例であった.また,2000〜2007年の
各年度末のゼロレートカーブ(後述)はいずれも右上がり,すなわち順イールドである.

\section{ゼロ(クーポン)・レート}

\subsection{最終利回りは使い難い?}

属性のうち,クーポンレートだけが異なる2つの債券の価格を考える.
\[
P_1 = \frac{F}{(1+Y_1)^N} + \sum_{i=1}^{N}\frac{FC_1}{(1+Y_1)^i},
\qquad
P_2 = \frac{F}{(1+Y_2)^N} + \sum_{i=1}^{N}\frac{FC_2}{(1+Y_2)^i}.
\]
通常,クーポンレート $C$ が違えば,価格 $P$ も最終利回り $Y$ も異なる.
すると,次のような疑問が生じる.

\begin{itemize}
\item 同一発行体の同時点のキャッシュフローを割り引くとき,クーポンレートが違う
債券では\emph{異なる}最終利回りで割り引くことになるのか?
\item 同一債券では,どの時点のキャッシュフローも\emph{同じ}利回りで割り引くことに
なるのか?
\end{itemize}

これらの疑問への一つの答え方は,「上記の式を割引の式と深く考えるな」,すなわち
これらは単にキャッシュフローと価格の関係(価格を利回りで言い換えたもの)と
みなせ,というものである.

\subsection{期間を代表する利回り:ゼロレート}

残存期間ごとに一意に決まる「期間を代表する利回り」を使って現在価値への割引を
行えば,話はスッキリする.では,期間を代表する利回りとして何を使えばよいか.
最終利回り $Y$ は,債券の条件(クーポンレートなど)が違うと値が変わる,という
問題が生じる.これを回避するのが次のゼロレートである.

\begin{definition}[ゼロレート]
\textbf{ゼロレート} (zero rate),\textbf{ゼロクーポンレート} (zero-coupon rate)
とは,割引債の最終利回りのことである.(“ゼロクーポン債=割引債”である.)
\end{definition}

残存期間を決めると,ゼロレートは一つに決まる(もし異なる値をとるならば,同満期の
割引債で価格が違うことになってしまう).したがって,
\begin{itemize}
\item ゼロレートは残存期間の代表的な金利として使える.
\item 異なる債券のキャッシュフローでも,同じ時点のキャッシュフローは,残存期間で
決まるゼロレートで割り引けばよい.
\end{itemize}
この考え方の転換は,講義では「ゼロレート革命」と呼ばれた.

\subsection{ゼロレートによる債券価格評価}

設定:年1回利払い,$N$ 年満期,額面 $F$ 円.$P$ を現在の価格,$C_i$ を $i$ 年目の
クーポン,$r_i$ を $i$ 年のゼロレートとする.ディスカウントファクター(割引関数)
$DF(i)$ は「$i$ 年後の1円の現価」であり,
\begin{equation}
DF(i) = \frac{1}{(1+r_i)^i} = (1+r_i)^{-i}
\label{eq:df-zero}
\end{equation}
と書ける.このとき債券価格は
\begin{equation}
P = \frac{C_1}{1+r_1} + \frac{C_2}{(1+r_2)^2} + \cdots + \frac{C_N}{(1+r_N)^N}
+ \frac{F}{(1+r_N)^N}
= \sum_{i=1}^{N} C_i\,DF(i) + F\,DF(N).
\label{eq:price-zero}
\end{equation}
なお,$DF(i)$ は $i$ 年満期の(額面1円の)割引債価格でもある.

\subsection{ゼロレートとYTMによる価格式の比較}

ゼロレートを用いた価格式と,最終利回りを用いた価格式を並べると,
\[
P = \sum_{i=1}^{N}\frac{C_i}{(1+r_i)^i} + \frac{F}{(1+r_N)^N}
\qquad\text{vs.}\qquad
P = \sum_{i=1}^{N}\frac{C_i}{(1+Y)^i} + \frac{F}{(1+Y)^N}.
\]
両者の違いに注意する.形式的には,右の式(最終利回りの定義式)は,利回りの
期間構造をフラットと仮定して CF を割り引いて現価を求めている形である.しかし,
これを「割引」とは考えないほうが,おそらくスッキリする(最終利回りは価格の
言い換えに過ぎない).

\subsection{連続複利における割引関数}

$DF(T) = D(T)$ を連続複利で表す.ゼロレートを $r_i$ として,$1/m$ 年複利で
$m\to\infty$ のケースを考えると,
\begin{equation}
D(T) = \lim_{m\to\infty}\Bigl(1+\frac{r_i}{m}\Bigr)^{-mT}
= \lim_{m\to\infty}\Biggl[\Bigl(1+\frac{1}{m/r_i}\Bigr)^{m/r_i}\Biggr]^{-r_i T}
= e^{-r_i T}.
\label{eq:df-cont}
\end{equation}
これが連続複利の(すなわち連続時間モデルにおける複利の)割引関数である.

\begin{remark}[イールドカーブの縦軸とスポット・レートという言葉]
イールドカーブ(金利期間構造)の縦軸を,債券の最終利回りでなく,ゼロレートを
使って表示することもよく行われる.特に理論の論文では,ゼロレートの方がよく
用いられる.一方,債券実務では今も最終利回りを使うことが多く,かつ取引は今も
単利である.

なお,実務では,現時点で取引されているゼロレートを\textbf{スポット・レート}
(spot rate) と呼ぶことがある.もともと「スポット」は直物(現時点で取引される
商品)を意味する.しかし,理論でスポットレートというと,少し別の意味(瞬間的な
短期金利という意味,本章付録参照)で使われることも多い.「スポット」は複数の
意味で使われる言葉なので,注意を要する.
\end{remark}

\section{フォワード・イールド}

本節の言葉の使い方はやや理論家寄りである(言葉の使い方はその人次第であり,
証券アナリストのテキストとは異なることもある).引き続き,安全確実な債券
($\approx$ 国債)を考える.

\subsection{イールド}

時刻を $t$ で表す.現在を $t=0$ として,$0 \le t \le T \le \tau$ とする.

\begin{definition}[イールド]
$Y(t,T)$:\textbf{イールド}.期間 $[t,T]$ の平均利回り.
\begin{itemize}
\item $t=0$ ならば,$Y(0,T)$ は確定値.
\item $t>0$ ならば,$Y(t,T)$ は確率変数(現時点 $t=0$ で値は確定しない).
\end{itemize}
\end{definition}

以下では,割引債による運用を考えて利回りについて記述する.割引債は株式と違い
キャッシュフローが確定的で考えやすいからである(株式の場合は,利回りと言っても
よいが,収益率と呼ぶ).定義より,$Y(0,T)$ が時刻 $t=0$ における期間 $T$ の
ゼロレートに相当する.

\subsection{運用と運用成果}

ある資産での運用を考える.時刻 $t$ における1単位あたりの価格を $p(t)$ とする.
期間 $[t,T]$ の運用成果を考える.投資額を $A$ 円とすると,
\begin{itemize}
\item $A / p(t)$:時刻 $t$ における購入量.
\item $A \times p(T)/p(t)$:時刻 $T$ における運用成果.
\item $p(T)/p(t)$:期間 $[t,T]$ の $(1+\text{利回り})$.
\end{itemize}

\begin{example}
期間 $[t,T]$ における,満期 $T$ の割引債(価格 $v(t,T)$)による(時刻 $t$ の)
1円の(時刻 $T$ における)運用成果は,
\[
1 \times \frac{v(T,T)}{v(t,T)} = \frac{1}{v(t,T)}.
\qquad(\because\ v(T,T)=1.)
\]
\end{example}

\subsection{フォワード・イールド}

\begin{definition}[フォワード・イールド]
時点 $t$ でみた,将来の期間 $[T,\tau]$($t \le T \le \tau$)の利回りを
\textbf{フォワード・イールド}といい,$f(t,T,\tau)$ と書く.
\end{definition}

定義から $Y(T,\tau) = f(T,T,\tau)$ である.$Y(T,\tau)$ は普通は時刻 $T$ に
ならないと確定しない.そんな利回りを $T$ 以前の時刻 $t$ で予想しても,その
予想値は確定値であろうはずがない.\emph{しかし},時刻 $t$ で取引されている
満期の異なる割引債を使うことで,この利回りを時刻 $t$ で確定させることもできる.
これが次のインプライド・フォワード・イールドである.

\subsection{現時点でフォワード・イールドは確定できる}

とりあえず $t=0$ として考える.いま取引されている(=スポットの)ゼロレートを
使うことで,将来のある期間の利回りを確定することができる.

\begin{example}[フォワード・イールドの確定]
ゼロレートが1年で5\%,2年で6\%とする.このとき,
\begin{itemize}
\item 満期1年の割引債(額面105万円)を空売りし,代金100万円を受け取る.
ここで $v(0,1) = DF(1) = 1/(1+0.05)$ なので,額面105万円の1年割引債の現在の
価格はちょうど $105 \times v(0,1) = 100$ 万円である.
\item その全額100万円で,満期2年の割引債(額面112.36万円)を買う.
$v(0,2) = DF(2) = 1/(1+0.06)^2$ なので,100万円で買える額面は
$100/v(0,2) = 100\times 1.06^2 = 112.36$ 万円である.
\end{itemize}
この合成ポジションのキャッシュフローは次のとおりである.
\begin{center}
\begin{tabular}{c|cc|cc|cc}
\toprule
時点(年) & \multicolumn{2}{c|}{1年割引債の売り} & \multicolumn{2}{c|}{2年割引債の買い}
& \multicolumn{2}{c}{合計}\\
 & 収入 & 支出 & 収入 & 支出 & 収入 & 支出\\
\midrule
0 & 100 & & & 100 & &\\
1 & & 105 & & & & 105\\
2 & & & 112.36 & & 112.36 &\\
\bottomrule
\end{tabular}
\end{center}
合計をみると,現時点のキャッシュフローはゼロで,「1年後に105万円を支払い,
2年後に112.36万円を受け取る」という取引になっている.すなわち,\textbf{この
取引で,1年後から2年後までの期間の利回りが現時点で確定}している.その利回りは
$112.36/105 - 1 \approx 7.01\%$ である.
\end{example}

\subsection{インプライド・フォワード・イールド}

\begin{definition}[インプライド・フォワード・イールド]
前項の方法で(市場で取引されている割引債・ゼロレートを使って)確定させた
フォワード・イールドを,\textbf{インプライド (implied)・フォワード・イールド}と
いう.
\end{definition}

この値は,1年物と2年物のゼロレートによって決まる値,つまり時刻 $t$ における
市場の状態によって決まる値である.

\begin{remark}
これを単にフォワード・イールドと呼ぶ人も,フォワード・レートと呼ぶ人も,
インプライド・フォワード・レートと呼ぶ人もいる.そのため,言葉の意味を確認
しながら使わないと危ない.
\end{remark}

インプライド・フォワード・イールドの話は一般化できる.記法として
\begin{itemize}
\item $Y(i,j,j+k)$:時点 $i$ における $j$ 期間から $(j+k)$ 期間までの利回り,
\item $Y(i,i+j)$:時点 $i$ における $j$ 期間の利回り(ゼロレート),
\end{itemize}
を用いる.基本となる考え方は次である.

\begin{keypoint}[裁定の考え方]
どういう期間区分で運用しても,($N$ 年後の)終価は同じ.少なくとも均衡状態では
そのようになっているはずである.(瞬間的にはズレるし,別の理由でズレが維持される
こともあるが.)
\end{keypoint}

例えば3年間の運用は,「3年ゼロレートでの一括運用」「1年目はゼロレート
$Y(0,1)$,残り2年はフォワード $Y(0,1,3)$」「1年ずつ $Y(0,0,1)$, $Y(0,1,2)$,
$Y(0,2,3)$ とつなぐ運用」のいずれでも同じ終価になる.

\subsection{インプライド・フォワード・イールドの公式(1年複利)}

\begin{example}[例1:2年間の投資]
1年目をゼロレート $Y(0,1)$ で,2年目をフォワード・イールド $Y(0,1,2)$ で運用した
成果と,2年ゼロレート $Y(0,2)$ で運用した成果は等しいので,
\[
\bigl(1+Y(0,1)\bigr)\bigl(1+Y(0,1,2)\bigr) = \bigl(1+Y(0,2)\bigr)^2
\qquad\therefore\quad
Y(0,1,2) = \frac{(1+Y(0,2))^2}{1+Y(0,1)} - 1.
\]
\end{example}

\begin{example}[例2:$j$ 年後から $j+k$ 年後の間の投資]
\[
\bigl(1+Y(0,j)\bigr)^{j}\bigl(1+Y(0,j,j+k)\bigr)^{k}
= \bigl(1+Y(0,j+k)\bigr)^{j+k}
\]
\[
\therefore\quad
Y(0,j,j+k) = \Biggl(\frac{(1+Y(0,j+k))^{j+k}}{(1+Y(0,j))^{j}}\Biggr)^{1/k} - 1.
\]
\end{example}

このように,\textbf{インプライド・フォワード・イールドはゼロレートで書ける}.
すなわち,「ゼロレートの期間構造」から「インプライド・フォワード・イールド
(の期間構造)」を求めることができる.

\begin{example}[例3:1年間の投資]
\[
Y(0,0,1) = Y(0,1).
\]
すなわち,最初の1期間のフォワード・イールドはゼロレートそのものである.
\end{example}

\begin{example}[例4:$N$ 年間の投資]
\[
\bigl(1+Y(0,0,1)\bigr)\bigl(1+Y(0,1,2)\bigr)\cdots\bigl(1+Y(0,N-1,N)\bigr)
= \bigl(1+Y(0,N)\bigr)^{N}
\]
\[
\therefore\quad
Y(0,N) = \sqrt[N]{\bigl(1+Y(0,0,1)\bigr)\bigl(1+Y(0,1,2)\bigr)\cdots
\bigl(1+Y(0,N-1,N)\bigr)} - 1.
\]
\end{example}

逆に,\textbf{ゼロレートはインプライド・フォワード・イールドで書ける}.
例4の形から,ゼロレートはインプライド・フォワード・イールドの期間平均
(幾何平均)的な値であることがわかる.この関係から,次が成り立つ.

\begin{keypoint}[カーブの形状の関係]
インプライド・フォワード・イールド・カーブが右上がり(右下がり,水平)ならば,
ゼロレート・カーブも右上がり(右下がり,水平)である.この関係は常に成立する.
しかし,\textbf{逆は常に成り立つとは限らない}(ゼロレート・カーブが右上がりでも,
フォワード・カーブが右上がりとは限らない).
\end{keypoint}

\subsection{インプライド・フォワード・イールドの公式(連続複利)}

\begin{example}[{例1:2年間(一般には $[j,j+k]$)の投資}]
連続複利では,運用成果の等式は指数の和になる:
\[
e^{j\,Y(0,j)}\,e^{k\,Y(0,j,j+k)} = e^{(j+k)\,Y(0,j+k)}
\qquad\therefore\quad
Y(0,j,j+k) = \frac{(j+k)\,Y(0,j+k) - j\,Y(0,j)}{k}.
\]
\end{example}

\begin{example}[例2:$N$ 年間の投資]
\[
e^{Y(0,0,1)}\,e^{Y(0,1,2)}\cdots e^{Y(0,N-1,N)} = e^{N\,Y(0,N)}
\qquad\therefore\quad
Y(0,N) = \frac{Y(0,0,1) + Y(0,1,2) + \cdots + Y(0,N-1,N)}{N}.
\]
\end{example}

連続複利の場合,\textbf{ゼロレートはインプライド・フォワード・イールドの期間
加重平均値}(単純平均)になる.

\subsection{割引債価格との関係}

\paragraph{(1) $t=0$ の場合.}
現在($t=0$)について考える.時刻 $t=0$ の満期 $T$,額面1円の割引債価格を
$v(0,T)$ と書く.ゼロレート $Y(0,T)$ との関係は,
\begin{align}
\text{1年複利:}\quad & v(0,T) = \bigl(1+Y(0,T)\bigr)^{-T},\notag\\
\text{半年複利:}\quad & v(0,T) = \bigl(1+Y(0,T)/2\bigr)^{-2T},\notag\\
\text{$1/m$年複利:}\quad & v(0,T) = \bigl(1+Y(0,T)/m\bigr)^{-mT},\notag\\
\text{連続複利:}\quad & v(0,T) = e^{-T\,Y(0,T)}.
\label{eq:v0T}
\end{align}
$v(0,T)$ は時刻 $t=0$(現在)における割引関数で,$t=0$ で既知である.

\paragraph{(2) インプライド・フォワード・イールドと割引債価格($t=0$).}
$Y(0,T,\tau)$ をインプライド・フォワード・イールド($t=0$ では既知)とすると,
\begin{align}
\text{1年複利:}\quad & Y(0,T,\tau)
= \Bigl(\frac{v(0,T)}{v(0,\tau)}\Bigr)^{1/(\tau-T)} - 1,\notag\\
\text{半年複利:}\quad & Y(0,T,\tau)
= 2\Bigl[\Bigl(\frac{v(0,T)}{v(0,\tau)}\Bigr)^{1/2(\tau-T)} - 1\Bigr],\notag\\
\text{$1/m$年複利:}\quad & Y(0,T,\tau)
= m\Bigl[\Bigl(\frac{v(0,T)}{v(0,\tau)}\Bigr)^{1/m(\tau-T)} - 1\Bigr],\notag\\
\text{連続複利:}\quad & Y(0,T,\tau)
= \frac{\log\bigl(v(0,T)/v(0,\tau)\bigr)}{\tau-T}.
\label{eq:fwd-v}
\end{align}
これらの導出は演習問題3-4である.導出の鍵は,$[0,T]$ 間の運用成果は
$\mathrm{price}(T)/\mathrm{price}(0)$ であり,満期 $T$ までの割引債 $v(0,T)$ の
運用成果は $1/v(0,T)$($\because v(T,T)=1$)であることである.

\begin{remark}
理論では,割引債価格を額面1円で考えて,割引関数とみなすのが普通である.実務では,
割引債でも債券なので,額面100円で価格を考える.
\end{remark}

\paragraph{(3)(4) $t>0$ の場合.}
ここまでは現在($t=0$)だけ考えたので $Y(0,\cdot)$ や $Y(0,\cdot,\cdot)$ と
したが,$t>0$ にしても考え方は同じである.ただし,変数はすべて確率変数になる.
時刻 $t$ の満期 $T$,額面1円の割引債価格を $v(t,T)$ と書くと,
\begin{align}
\text{1年複利:}\quad & v(t,T) = \bigl(1+Y(t,T)\bigr)^{-(T-t)},\notag\\
\text{半年複利:}\quad & v(t,T) = \bigl(1+Y(t,T)/2\bigr)^{-2(T-t)},\notag\\
\text{$1/m$年複利:}\quad & v(t,T) = \bigl(1+Y(t,T)/m\bigr)^{-m(T-t)},\notag\\
\text{連続複利:}\quad & v(t,T) = e^{-(T-t)\,Y(t,T)},
\end{align}
であり($v(t,T)$, $t>0$ は時刻 $t$ の割引関数で,$t=0$(現在)では確率変数),
インプライド・フォワード・イールド $Y(t,T,\tau)$(確率変数)についても
\begin{align}
\text{1年複利:}\quad & Y(t,T,\tau)
= \Bigl(\frac{v(t,T)}{v(t,\tau)}\Bigr)^{1/(\tau-T)} - 1,\notag\\
\text{半年複利:}\quad & Y(t,T,\tau)
= 2\Bigl[\Bigl(\frac{v(t,T)}{v(t,\tau)}\Bigr)^{1/2(\tau-T)} - 1\Bigr],\notag\\
\text{$1/m$年複利:}\quad & Y(t,T,\tau)
= m\Bigl[\Bigl(\frac{v(t,T)}{v(t,\tau)}\Bigr)^{1/m(\tau-T)} - 1\Bigr],\notag\\
\text{連続複利:}\quad & Y(t,T,\tau)
= \frac{\log\bigl(v(t,T)/v(t,\tau)\bigr)}{\tau-T}
\end{align}
が成り立つ.

\section*{付録:スポット・レートとフォワード・レート}

デリバティブの実務では,これらが確率変動するモデルを使い,金利期間構造の変動を
議論する.

\subsection*{(瞬間的な)フォワードレート}

フォワード・イールド $f(t,T,\tau)$ の定義で,$\tau = T + dT$ として $dT\to 0$ の
極限を考えたものが,\textbf{(瞬間的な)フォワードレート}である:
\begin{equation}
f(t,T) = \lim_{\Delta T\to 0} Y(t,T,T+\Delta T).
\end{equation}
フォワード・レートは,時刻 $t$ でみた,$[T, T+dT]$ 間の一瞬の利回りである.

$t=0$ の場合を連続複利の式\eqref{eq:fwd-v}で考えると,
\begin{equation}
f(0,t) = \lim_{\Delta t\to 0} Y(0,t,t+\Delta t)
= \lim_{\Delta t\to 0}
\frac{(t+\Delta t)\,Y(0,t+\Delta t) - t\,Y(0,t)}{\Delta t}
= \frac{\dd\,[\,t\,Y(0,t)\,]}{\dd t}.
\end{equation}
これより,
\begin{equation}
\int_0^t f(0,s)\,\dd s = t\,Y(0,t)
\qquad\therefore\quad
Y(0,t) = \frac{1}{t}\int_0^t f(0,s)\,\dd s.
\end{equation}

同様に,$t>0$,つまり将来においても,満期変数($T$)についての微分として
\begin{equation}
f(t,T) = \frac{\dd\,[\,(T-t)\,Y(t,T)\,]}{\dd T}
\end{equation}
となるので,
\begin{equation}
\int_t^T f(t,s)\,\dd s = (T-t)\,Y(t,T)
\qquad\therefore\quad
Y(t,T) = \frac{1}{T-t}\int_t^T f(t,s)\,\dd s.
\end{equation}

このように,現在および将来のゼロレートはフォワードレートで記述できる.
したがって,\textbf{フォワードレート $f(t,T)$ の変動過程を上手にモデル化すれば,
金利期間構造の全体の変動を上手に表現できる}.これが現在の確率金利モデルの
基本的な考え方である.

\subsection*{(瞬間的な)スポットレート}

(瞬間的な)フォワードレート $f(t,T)$ の $T \to t$ の極限値を,
\textbf{(瞬間的な)スポットレート}という:
\begin{equation}
r(t) = \lim_{T\to t} f(t,T).
\end{equation}
スポット・レートは,時刻 $t$ でみた,$[t, t+dt]$ の一瞬の利回りである.
スポットレート $r(t)$ の変動過程を上手にモデル化しても,金利期間構造の変動を
上手に表現できる.デリバティブの価格付けでは,$r(t)$ や $f(t,T)$ を確率過程と
して扱ったモデルがよく使われる.

% ---------------- 演習問題 ----------------
\section{演習問題}

\begin{exercise}[問題3-1:債券価格と最終利回り]
額面100円,満期3年,クーポンレート2\%,半年払い($m=2$)の債券の,将来
キャッシュフローと現在価値を求めなさい.ただし,半年複利の最終利回りは3\%と
する.
\end{exercise}

\begin{answerbox}
\textbf{考え方.}まずキャッシュフローを書き出す.クーポンは
$C_i = K\times F\times(t_i - t_{i-1}) = 2\%\times 100\times 0.5 = 1$ 円
(半年ごと)であり,満期3年には額面100円が加わって101円となる.次に,各時点の
1円の現価(ディスカウントファクター)$DF(t) = (1+0.03/2)^{-2t}$ を掛けて
合計する.

\medskip
\textbf{計算.}
\begin{center}
\begin{tabular}{cccc}
\toprule
年 & キャッシュフロー & 1円の現価 & キャッシュフローの現価\\
\midrule
0.5 & 1 & 0.9852 & 0.9852\\
1.0 & 1 & 0.9707 & 0.9707\\
1.5 & 1 & 0.9563 & 0.9563\\
2.0 & 1 & 0.9422 & 0.9422\\
2.5 & 1 & 0.9283 & 0.9283\\
3.0 & 101 & 0.9145 & 92.3688\\
\midrule
\multicolumn{3}{r}{債券価格} & 97.1514\\
\bottomrule
\end{tabular}
\end{center}
よって現在価値(債券価格)は約 \textbf{97.15円}である.

\medskip
\textbf{ポイント.}クーポンレート(2\%)が最終利回り(3\%)より低いので,価格は
額面100円を下回る(アンダーパー).「最終利回り $=$ クーポンレート $\iff$
価格 $=$ 額面」という関係とあわせて確認しておくとよい.
\end{answerbox}

\begin{exercise}[問題3-2:インプライド・フォワード・イールドとゼロレート(離散時間モデル)]
次の表の空欄のインプライド・フォワード・イールドとゼロレートの値を求めなさい.
ただし,1年を一期間とする離散時間モデルで考える.
\begin{center}
\small
\begin{tabular}{c|cc|cc|cc}
\toprule
 & \multicolumn{2}{c|}{ケースA} & \multicolumn{2}{c|}{ケースB}
 & \multicolumn{2}{c}{ケースC}\\
年 & フォワード & ゼロ & フォワード & ゼロ & フォワード & ゼロ\\
\midrule
1 & 7.0\% & 7.000\% & 5.0\% & --- & --- & 9.0\%\\
2 & 7.0\% & --- & 6.0\% & --- & --- & 8.5\%\\
3 & 7.0\% & --- & 7.0\% & --- & --- & 8.0\%\\
4 & 7.0\% & --- & 8.0\% & --- & --- & 7.5\%\\
5 & 7.0\% & --- & 9.0\% & --- & --- & 7.0\%\\
\bottomrule
\end{tabular}
\end{center}
※ 表の「フォワードレート」は「インプライド・フォワード・イールド」に読み替える.
表中のフォワード・イールドは,それぞれ運用期間1年間のインプライド・フォワード・
イールドの意味である.例えば,ケースBの5年の9.0\%は $f(0,4,5)$,つまり現時点で
みた4年から5年の1年間の投資期間の利回りであり,4年の8.0\%は $f(0,3,4)$ である.
\end{exercise}

\begin{answerbox}
\textbf{考え方.}$n$ 年のゼロレート $r_n$ と1年刻みのフォワード・イールド
$f_j = f(0,j-1,j)$ の間には,「どの期間区分で運用しても $n$ 年後の終価は同じ」
という関係
\[
(1+r_n)^n = (1+f_1)(1+f_2)\cdots(1+f_n)
\]
が成り立つ.フォワードからゼロを求めるにはこの式をそのまま使い,ゼロから
フォワードを求めるには
\[
1+f_n = \frac{(1+r_n)^n}{(1+r_{n-1})^{n-1}}
\]
を使う.

\medskip
\textbf{計算結果.}
\begin{center}
\small
\begin{tabular}{c|cc|cc|cc}
\toprule
 & \multicolumn{2}{c|}{ケースA} & \multicolumn{2}{c|}{ケースB}
 & \multicolumn{2}{c}{ケースC}\\
年 & フォワード & ゼロ & フォワード & ゼロ & フォワード & ゼロ\\
\midrule
1 & 7.0\% & 7.000\% & 5.0\% & 5.000\% & 9.000\% & 9.0\%\\
2 & 7.0\% & 7.000\% & 6.0\% & 5.499\% & 8.002\% & 8.5\%\\
3 & 7.0\% & 7.000\% & 7.0\% & 5.997\% & 7.007\% & 8.0\%\\
4 & 7.0\% & 7.000\% & 8.0\% & 6.494\% & 6.014\% & 7.5\%\\
5 & 7.0\% & 7.000\% & 9.0\% & 6.991\% & 5.023\% & 7.0\%\\
\bottomrule
\end{tabular}
\end{center}
例えばケースBの2年ゼロレートは
$r_2 = \sqrt{1.05\times 1.06} - 1 = 5.499\%$,ケースCの2年フォワードは
$f_2 = 1.085^2/1.09 - 1 = 8.002\%$ である.

\medskip
\textbf{ポイント.}
\begin{itemize}
\item ケースA:フォワードが常に7\%なら,ゼロレートも常に7\%(フラット).
\item ケースB:フォワードが右上がりのとき,ゼロレートも右上がりになるが,
ゼロレートはフォワードの「平均」なので,伸びはフォワードより緩やか.
\item ケースC:ゼロレートが右下がりのとき,フォワードはそれより急な右下がりに
なる.
\end{itemize}
ゼロレートはインプライド・フォワード・イールドの平均値的なものである.しかし,
離散時間モデルでは(幾何平均のため)微妙にズレる.連続時間モデルでは,まさに
期間平均値になる.
\end{answerbox}

\begin{exercise}[問題3-3:ゼロレートと債券価格]
額面100円,残存期間3年,クーポンレート6\%(年1回払い)の満期一括償還の利付債を
考える.
\begin{enumerate}
\item ゼロレートが1年で5\%,2年で6\%,3年で7\%のとき,利付債の価格と最終利回りを
求めなさい.
\item ゼロレートが1年で7\%,2年で6\%,3年で5\%のとき,利付債の価格と最終利回りを
求めなさい.
\end{enumerate}
ただし,1年を一期間とする離散時間モデルで考えること.
\end{exercise}

\begin{answerbox}
\textbf{考え方.}価格は,各時点のキャッシュフロー(1・2年目:6円,3年目:106円)を
それぞれの残存期間のゼロレートで割り引いて合計する(式\eqref{eq:price-zero}).
最終利回りは,得られた価格 $P$ に対して
$P = \sum_{i=1}^{3} CF_i/(1+Y)^i$ を満たす $Y$ を数値的に求める.

\medskip
\textbf{1.(順イールドのケース)}
\begin{center}
\small
\begin{tabular}{ccccc}
\toprule
年 & Cash Flow & Zero Rate & Discount Function & Present Value\\
\midrule
1 & 6.00 & 5.000\% & 0.952381 & 5.71429\\
2 & 6.00 & 6.000\% & 0.889996 & 5.33998\\
3 & 106.00 & 7.000\% & 0.816298 & 86.52757\\
\midrule
\multicolumn{4}{r}{PRICE} & 97.58184\\
\multicolumn{4}{r}{YTM} & 6.920\%\\
\bottomrule
\end{tabular}
\end{center}

\textbf{2.(逆イールドのケース)}
\begin{center}
\small
\begin{tabular}{ccccc}
\toprule
年 & Cash Flow & Zero Rate & Discount Function & Present Value\\
\midrule
1 & 6.00 & 7.000\% & 0.934579 & 5.60748\\
2 & 6.00 & 6.000\% & 0.889996 & 5.33998\\
3 & 106.00 & 5.000\% & 0.863838 & 91.56679\\
\midrule
\multicolumn{4}{r}{PRICE} & 102.51424\\
\multicolumn{4}{r}{YTM} & 5.075\%\\
\bottomrule
\end{tabular}
\end{center}

\medskip
\textbf{ポイント.}どちらのケースもゼロレートは 5\%,6\%,7\% の同じ3つの値から
なるが,並び順が違うだけで価格も最終利回りも大きく異なる.1.では YTM
$=6.920\%$ と3年ゼロレート(7\%)に近く,2.では YTM $=5.075\%$ と3年ゼロレート
(5\%)に近い.これは,この債券のキャッシュフローの大部分(106円)が3年目に
あるためである.すなわち,\textbf{YTM は,キャッシュフローが大きな時点の
ゼロレートを強く反映}する.
\end{answerbox}

\begin{exercise}[問題3-4:割引債価格との関係]
\begin{enumerate}
\item[(1)] 割引債価格とゼロレートの関係式\eqref{eq:v0T}を導出しなさい.
\item[(2)] インプライド・フォワード・イールドと割引債価格の関係式
\eqref{eq:fwd-v}を導出しなさい.
\end{enumerate}
\end{exercise}

\begin{answerbox}
\textbf{(1) の解答.}
割引債価格 $v(0,T)$ は,$t=0$(現在)における割引関数 $DF(T)$ とみなせる
(額面1円の割引債の価格は「$T$ 年後の1円の現価」そのものである).よって,
離散時間モデルの複利表示の式は,ディスカウントファクターの式
\eqref{eq:df-zero}(を $1/m$ 年複利に一般化したもの)
\[
v(0,T) = DF(T) = \Bigl(1+\frac{Y(0,T)}{m}\Bigr)^{-mT}
\]
から得られる.連続時間モデルの連続複利表示の式は,連続複利の割引関数の式
\eqref{eq:df-cont}から $v(0,T) = e^{-T\,Y(0,T)}$ として得られる.

\medskip
\textbf{(2) の解答.}
代表として半年複利の式(2番目の式)を示す.0年から $T$ 年まで割引債で運用し,
$T$ 年から $\tau$ 年までインプライド・フォワード・イールドで運用した成果は
\[
\frac{1}{v(0,T)}\times\Bigl(1+\frac{Y(0,T,\tau)}{2}\Bigr)^{2(\tau-T)}.
\]
一方,0年から $\tau$ 年まで割引債で運用したときの成果は $1/v(0,\tau)$ である.
均衡ではこれらは等しいので,
\[
\frac{1}{v(0,T)}\Bigl(1+\frac{Y(0,T,\tau)}{2}\Bigr)^{2(\tau-T)}
= \frac{1}{v(0,\tau)}
\]
が成り立ち,これを $Y(0,T,\tau)$ について解けば
\[
Y(0,T,\tau) = 2\Bigl[\Bigl(\frac{v(0,T)}{v(0,\tau)}\Bigr)^{1/2(\tau-T)} - 1\Bigr]
\]
を得る.1年複利・$1/m$ 年複利の式(1番目と3番目の式)も同様に得られる.

連続複利の式(4番目の式)も同様である.0年から $T$ 年まで割引債で運用し,
$T$ 年から $\tau$ 年までインプライド・フォワード・イールドで運用した成果は
\[
\frac{1}{v(0,T)}\times e^{(\tau-T)\,Y(0,T,\tau)}
\]
であり,これが $1/v(0,\tau)$ に等しいことから
\[
\frac{1}{v(0,T)}\,e^{(\tau-T)\,Y(0,T,\tau)} = \frac{1}{v(0,\tau)}
\qquad\therefore\quad
Y(0,T,\tau) = \frac{\log\bigl(v(0,T)/v(0,\tau)\bigr)}{\tau-T}.
\]

\medskip
\textbf{ポイント.}導出の本質は「運用成果の一致(裁定が働けば,どの経路で運用
しても同じ期間の確定運用成果は等しい)」である.割引債の運用成果が $1/v(0,T)$ と
書けること($v(T,T)=1$)さえ押さえれば,あとは機械的に式変形できる.
\end{answerbox}
